\subsection{局部凸空间的例子}

\begin{enumerate}
\def\labelenumi{\arabic{enumi})}
\item
  空间 \(R^{n}\) 是局部凸的，因为以 0 为中心的开立方体都是凸集 (II, 第 9 页, prop. 6)。因此，这对所有有限维实拓扑向量空间也成立；事实上，只要 E 是 Hausdorff 的，由上述事实及 I, § 2.3, th. 2 即得；若不然，与 E 相伴的 Hausdorff 空间 F 是有限维的，因而是局部凸的，而 F 中 0 的凸邻域在典范映射 \(E \rightarrow F\) 下的逆像是凸的，并且构成 E 中 0 的一个基本邻域系。
\item
  设 E 是 R 上的向量空间，B 是 E 中所有吸收、对称且凸的子集所成的族。由 II, 第 23 页, prop. 1 可见，B 是 E 上某个局部凸拓扑 \(T_{\omega}\) 的 0 的基本邻域系，这个拓扑是 E 上所有局部凸拓扑中最细的。这个拓扑是 Hausdorff 的；因为设 \(x \neq 0\) 为 E 的任意一点；存在 E 的一个基 \((i_{\mathfrak{t}})_{\mathfrak{t} \in \mathbf{I}}\) 以及某个 \(\alpha \in \mathbf{I}\) 使 \(e_{\alpha} = x\) ；点 \(y = \sum_{\mathfrak{t}} y_{\mathfrak{t}} e_{\mathfrak{t}}\)（满足 \(|y_{\alpha}| < 1\)）所成的集是吸收、对称且凸的。它不含 x。
\end{enumerate}

由 II, 第 24 页, 推论可知，\(T_{\omega}\) 也是由 E 上全体半范数所成的族定义的拓扑，因此每个半范数在 \(T_{\omega}\) 下都连续。

特别地，若 u 是 E 到任一局部凸空间 F 中的线性映射，则 F 中 0 的每个凸邻域在 u 下的逆像是 E 中的吸收凸集；因此它是 \(T_{\omega}\) 下 0 的一个邻域，从而 u 对 \(T_{\omega}\) 连续。

对 E 中的凸集 C，我们称点 \(a \in C\) 为 C 的代数内点，如果对每条包含 a 的直线 D，交 \(D \cap C\) 包含一个包含 a 的开线段；换言之，\(-a + C\) 是吸收的。E 中集 A 的点 a 对 \(T_{\omega}\) 是 A 的内点，必须且只需存在凸集 C 使 \(a \in C \subset A\) ，并且 a 是 C 的代数内点。

更一般地，设 V 是 E 中的一个仿射线性流形，C 是含于 V 中的凸集；点 \(a \in C\) 称为 C 相对 V 的代数内点，如果在向量子空间 \(V_{0} = -a + V\) 中，点 0 是集 \(C_{0} = -a + C\) 的代数内点。

当 E 有限维时，拓扑 \(T_{\omega}\) 就是 E 上的典范拓扑 (I, 第 13 页, th. 2)；这表明 E 中凸集 C 的每个代数内点对典范拓扑都是 C 的内点（参见 II, 第 74 页, exerc. 5）。

\begin{enumerate}
\def\labelenumi{\arabic{enumi})}
\setcounter{enumi}{2}
\tightlist
\item
  设 A 是 R 上向量空间 E 中的对称凸集。由 A 生成的向量子空间 F 也是由 A 生成的凸锥，因为 -A = A；这个集是形如 \(\lambda x\) （其中 \(x \in A\) 且 \(\lambda \in R\) ）的元所成的集；集 A 在 F 中是吸收的，而诸集 \(\lambda A\) （其中 \(\lambda > 0\) ）构成 F 上某个局部凸拓扑（称为由 A 定义的拓扑）的 0 的一个基本邻域系，该拓扑由半范数 \(p_{A}\) 即 A 的度规定义 (II, 第 20 页, prop. 22)；我们把赋予 F 这个半范数所得到的局部凸空间记作 \(E_{A}\) 。空间 \(E_{A}\) 是 Hausdorff 的，必须且只需 \(p_{A}\) 是一个范数，或者等价地，A 不包含任何直线。若 B 是 E 中第二个对称凸集且 \(A \subset B\) ，则显然 \(E_{A} \subset E_{B}\) ，且 \(E_{A}\) 到 \(E_{B}\) 中的典范单射对分别由 A 与 B 定义的拓扑连续。此外，若 f 是 E 到实向量空间 \(E'\) 中的线性映射，则 \(f(A)\) 在 \(E'\) 中是凸的且对称的，而 f 是 \(E_{A}\) 到 \(E'_{f(A)}\) 上的连续线性映射。
\end{enumerate}

最后注意，若 E 带有与其向量空间结构相容的拓扑 T，且 V 是 T 下 0 的一个对称凸邻域，则由 V 生成的向量空间与 E 相同，因为 V 是吸收的，而且 E 到 \(E_{V}\) 中的恒等映射是连续的。

